\subsection{Uniform tame coefficients: a root-stack argument}

\begin{proposition}[F8: uniform tame coefficient bound; Proved]
Let $A$ be a complete normal two-dimensional local domain of
characteristic $p>0$ with algebraically closed residue field $k$.
Put $U=\Spec A\setminus\{\mathfrak m\}$, and fix an SNC resolution
$X\to\Spec A$ with reduced exceptional divisor
$D=\bigcup_{i=1}^s C_i$. There is a constant $C_X$ such that
\[
 \dim_{\mathbf F_p}H^1(U,M)\leq C_X\operatorname{rank}M
 \tag{F8.0}
\]
for every finite locally constant $\mathbf F_p$-sheaf $M$ on $U$
whose monodromy is tamely ramified along $D$.
The finite monodromy group need not be abelian or of order prime to $p$.
\end{proposition}

\paragraph{Status.}
The complete proof passed an independent root audit, including the
primary sources and nearby-cycle formulas; BCE separately checked
the root-gerbe degree, coarse-descent, and character-count steps.
This is a proved result of the present work, without a historical
priority claim. The uniformity in root orders is part of the proof.

\begin{proof}
Choose a coefficient field $k\hookrightarrow A$. Cohen structure
expresses $A$ as a quotient of a power-series ring over perfect $k$,
which is finite over its Frobenius image; hence $A$ and the
essentially finite-type chart rings below are $F$-finite. If $A$ is regular, first blow up its closed point,
so $D$ is nonempty. Write
\[
 B_{ij}=C_i\cdot C_j,\qquad g_i=g(C_i),\qquad
 \nu_i=\#\{\text{crossings on }C_i\}.
\]
Choose positive integers $c_i$ so that
$J=\mathcal O_X(-Z)$, $Z=\sum_i c_iC_i$, has positive degree
\[
 d_i=-\sum_jc_j B_{ij}>0
\]
on every component. Such a cycle exists by the negative definiteness
used in Suzuki's proof. Set
\[
 L_i=\sum_jc_j|B_{ij}|,\qquad
 K_i=\max\{0,-B_{ii},2g_i-2+\nu_i\}.
 \tag{F8.1}
\]
These constants depend only on the chosen resolution.

\paragraph{1. Root stacks, extension, and degree conventions.}
Choose an integer $N$ prime to $p$ divisible by the tame inertia
orders of the finite monodromy of $M$ along all $C_i$. Let
\[
 \pi:\mathfrak X_N=\sqrt[N]{(X,C_1)}\mathbin{\times_X}\cdots
                 \mathbin{\times_X}\sqrt[N]{(X,C_s)}\longrightarrow X.
\]
Locally at a crossing this has the quotient chart
\[
 \left[\Spec R[u,v]/(u^N-x,v^N-y)\,/\,\mu_N^2\right],
\]
with the analogous one-variable chart away from crossings.
These charts are regular, and the stack is a regular tame
Deligne--Mumford stack, proper over $A$, with coarse space $X$.
Let $\mathfrak C_i$ be its reduced divisor above $C_i$. Then
$N\mathfrak C_i=\pi^*C_i$.

The sheaf $M$ extends to a finite locally constant sheaf on
$\mathfrak X_N$. Indeed the displayed root charts kill its
codimension-one tame inertia by the DVR form of Abhyankar's lemma.
Normalization on the regular charts is then \'etale by purity;
the canonical normalizations and their descent data agree on
chart overlaps. This proves the asserted extension, using only
Abhyankar's lemma and regular-scheme purity
\cite[tags 0BRM and 0BMB]{StacksAbhyankar,StacksBranchPurity}.
The same argument applies if $M$ has wild monodromy away from
these inertia groups.

The component $\mathfrak C_i$ is a $\mu_N$-gerbe over the root
curve obtained from $C_i$ by inserting $N$th roots at its $\nu_i$
crossing points. Degrees on $\mathfrak C_i$ will be normalized
so that the pullback of a line bundle on $C_i$ has its ordinary
degree. Concretely every line bundle has its $N$th power descending
to $C_i$, and we divide that descended degree by $N$. This convention
is additive and gives
\[
 \deg\bigl(\mathcal O(\mathfrak C_j)|_{\mathfrak C_i}\bigr)
        =B_{ij}/N,
 \qquad \deg(J|_{\mathfrak C_i})=d_i.
 \tag{F8.2}
\]
For vector bundles use the degree of the determinant. Nonzero
maps of line bundles do not increase degree, since their zeros
are effective divisors of nonnegative rational degree.
The degree multiplies by $p$ under absolute Frobenius, since
Frobenius pulls a line bundle back to its $p$th power.
The local charts and degree convention agree with the usual
root-stack/parabolic description
\cite[\S\S2.3--2.4]{Cadman2007}
\cite[Theorems 4--5, pp.~12,17]{Borne2007}.

Set $V=M\otimes_{\mathbf F_p}\mathcal O_{\mathfrak X_N}$, of rank $r$.
It has a unit Frobenius structure $F^*V\simeq V$ and is trivialized
by a finite \'etale cover. Consequently $V|_{\mathfrak C_i}$ is
semistable of degree zero: pull a subbundle back to that cover,
take its determinant inside an exterior power of a trivial bundle,
and use the preceding degree inequality. Under a finite representable
cover the normalized degrees scale by a positive rational factor:
ordinary stack degree scales by the degree of the cover, while
normalizing multiplies by the generic stabilizer order. Thus degree
signs are preserved. Its determinant has degree zero because it
becomes trivial on the cover. No restriction
on the degree of the trivializing cover is used here. For root curves with trivial generic stabilizer this is also
the finite-bundle semistability assertion of
\cite[Proposition 6, p.~22]{Borne2007}.

\paragraph{2. Coarse cohomology and a useful curve estimate.}
For a vector bundle $W$ on $\mathfrak C_i$, decompose it by the
characters of the generic $\mu_N$. All nontrivial generic-character
summands have zero coherent cohomology: taking invariants along
that gerbe is exact and kills them. The character-zero summand
$W_0$ descends to the root curve of $C_i$ at the crossings.
Its pushforward $E_0$ to $C_i$ is a vector bundle of the same rank.
On a local root chart a graded basis gives
\[
 0\leq \deg W_0-\deg E_0\leq\nu_i\operatorname{rank}W_0.
 \tag{F8.3}
\]
For clarity, adjunction from the pullback of $E_0$ is a generically
invertible map to $W_0$. In a character basis at a stacky point,
its matrix has diagonal entries $z^{a_j}$ with $0\leq a_j<N$.
The determinant degree lost in pushforward is $\sum_j a_j/N$.
Summing over the crossings proves (F8.3). Graded bases are obtained
by lifting a character basis modulo the maximal ideal and using
Nakayama's lemma; averaging makes the lifts equivariant.

Thus if $W_0$ is semistable of negative slope $t$, then it has no
sections and ordinary Riemann--Roch for $E_0$ gives
\[
 h^1(\mathfrak C_i,W_0)
 \leq \operatorname{rank}(W_0)(g_i-1-t+\nu_i).
 \tag{F8.4}
\]
If its slope is zero, evaluation at a nonstacky point is injective
on sections, since a section vanishing there gives a positive-degree
line subbundle after saturation. Hence
\[
 h^1(\mathfrak C_i,W_0)
 \leq \operatorname{rank}(W_0)(g_i+\nu_i).
 \tag{F8.5}
\]
These statements also apply to an isotypic summand tensored by a
line bundle that makes its generic character zero. Each isotypic
summand of a degree-zero semistable bundle is again semistable of
degree zero: it is a direct summand, every summand has degree at
most zero, and their degrees sum to zero.

Coarse pushforward on these tame stacks is exact on coherent
sheaves and on sheaves of $\mathbf F_p$-vector spaces. Locally it
is invariants for a finite group of order prime to $p$, followed
by finite affine pushforward; the invariants functor is exact by
averaging. For the \'etale assertion, take a strict henselian
neighborhood on the coarse space: its finite root cover is a finite
union of strictly henselian local schemes, whose \'etale global
section functors are exact. The quotient atlas then computes the
pushforward by finite-group cohomology, which vanishes in positive
degrees. The same description commutes with restriction to a closed
fibre. More explicitly, write $j:U\hookrightarrow\mathfrak X_N$ and
$j_X:U\hookrightarrow X$. Exact coarse pushforward gives
$R\pi_*Rj_*M=Rj_{X*}M$, and its closed-fibre compatibility gives
$i_X^*R\pi_*Rj_*M=R\pi_{D*}i^*Rj_*M$. These identities apply
to the bounded-below torsion complex $Rj_*M$ degree by degree.
Proper base change for the proper scheme $X$ over the strictly
henselian base $\Spec A$ therefore identifies
$H^1(U,M)$ with $\mathbb H^1(\mathfrak D_N,i^*Rj_*M)$.
Its spectral sequence yields the required nearby-cycle five-term
sequence. Coherent cohomology is computed by the same exact
coarse pushforward and the preceding coarse-space estimates.

\paragraph{3. A Frobenius-cokernel slope bound independent of $N$.}
Let
\[
 B_N=(F_*\mathcal O_{\mathfrak X_N})/\mathcal O_{\mathfrak X_N}.
\]
It is locally free: absolute Frobenius induces the automorphism
$\zeta\mapsto\zeta^p$ on each $\mu_N$ stabilizer, so it is
representable and is finite flat on the regular root charts. Its
unit section extends locally to a basis.
The pullback $F^*B_N$ is canonically isomorphic, as a module, to
the augmentation ideal in $F^*F_*\mathcal O$: the unit provides a
section of the augmentation. The powers of that ideal give a
filtration with locally free graded pieces
\[
 T^a(\Omega^1_{\mathfrak X_N/k}),\qquad 1\leq a\leq2(p-1),
 \tag{F8.6}
\]
where $T^a$ is the degree-$a$ truncated symmetric power.
This is Sun's canonical filtration
\cite[Theorem 3.7, pp.~12--13; local proof pp.~6--8]{Sun2008}.
In the present complete-local setting it can also be checked
directly after faithfully flat completion of a regular root chart:
with parameters $u,v$ and coefficient field $k$, the algebra is
\[
 k[[u,v]]\otimes_{k[[u^p,v^p]]}k[[u,v]]
 \simeq k[[u,v]][\alpha,\beta]/(\alpha^p,\beta^p),
\]
with augmentation ideal $(\alpha,\beta)$. Its graded generators
are $du,dv$. This calculation is invariant under changes of
parameters and therefore descends. All modules and maps involved are finite. Every closed point on
the proper exceptional divisor has residue field $k$, so its
completed regular chart is $k[[u,v]]$. Faithfully flat completion
checks these identities at every such point; any nonzero coherent
error supported on that divisor would have a closed point in its
support. Thus the calculation suffices on a neighborhood of the
exceptional divisor, which is all that the proof uses.

The conormal sequence on $\mathfrak C_i$ is an extension of two
line bundles, of degrees
\[
 -B_{ii}/N,\qquad
 2g_i-2+\nu_i(1-1/N).
 \tag{F8.7}
\]
The second formula follows from $dx=Nz^{N-1}dz$ at each root point;
the generic gerbe is \'etale and contributes no cotangent direction.
Both numbers are at most $K_i$. Each $T^a(\Omega^1)|_{\mathfrak C_i}$
has a filtration by the characters of this two-line flag: its
weights are $u,v$ with $0\leq u,v\leq p-1$ and $u+v=a$.
Thus its line-bundle layers have degrees at most $2(p-1)K_i$.
Tensoring with the semistable degree-zero bundle $V|_{\mathfrak C_i}$
shows that every subbundle of
$F^*(V\otimes B_N)|_{\mathfrak C_i}$ has slope at most
$2(p-1)K_i$. Here we used $F^*V\simeq V$ and the filtration (F8.6).
The assertion about subbundles follows by intersecting with the
filtration and adding degrees and ranks; each graded subbundle has
the indicated slope bound.

Pulling any subbundle of $(V\otimes B_N)|_{\mathfrak C_i}$ back by
Frobenius and dividing degrees by $p$ therefore proves the uniform
upper slope bound
\[
 \mu(\text{any subbundle of }(V\otimes B_N)|_{\mathfrak C_i})
       \leq b_i:=2(p-1)K_i/p.
 \tag{F8.8}
\]
This is independent both of $N$ and of $r$.

\paragraph{4. Nearby-cycle filtration and a uniform cutoff.}
Write $\mathfrak Z_N=\pi^*Z=\sum_iNc_i\mathfrak C_i$ and continue
to denote its ideal by $J$. The sheaf $M$ is locally constant on
$\mathfrak X_N$, so the Artin--Schreier calculation of F5/F7 applies
on its \'etale charts and descends. In particular, with
$\Psi=i^*j_*$ for the reduced exceptional stack,
\[
 0\to J^{-m}V/V\xrightarrow{F_V-1}J^{-pm}V/V
       \to U_mR^1\Psi M\to0,
 \tag{F8.9}
\]
and, for $m\geq2$,
\[
 \operatorname{gr}_mR^1\Psi M
  \simeq B_N|_{\mathfrak Z_N}\otimes V|_{\mathfrak Z_N}
                          \otimes J^{-m}|_{\mathfrak Z_N}.
 \tag{F8.10}
\]
The projection formula proves the latter identity, as in Suzuki's
Proposition 3.1; $F^{-1}\mathfrak Z_N=p\mathfrak Z_N$.

Filter $\mathcal O_{\mathfrak Z_N}$ by removing one irreducible
component at a time. A line-bundle layer $P$ on $\mathfrak C_i$
has degree bounded above by $L_i$, by (F8.2): all cycle coefficients
are between $0$ and $Nc_j$. The corresponding layer of (F8.10)
therefore has every subbundle of slope at most
\[
 b_i+L_i-md_i.
\]
Choose an integer $m_0\geq2$ strictly greater than
$\max_i(b_i+L_i)/d_i$. The last expression is negative for every
$m\geq m_0$, so there are no sections on any layer and hence
\[
 H^0(\operatorname{gr}_mR^1\Psi M)=0\quad(m\geq m_0).
 \tag{F8.11}
\]
The number of layers grows with $N$, but their common vanishing
bound does not. Quasi-compactness and the increasing union of the
pole lattices give
$H^0(R^1\Psi M)=H^0(U_{m_0}R^1\Psi M)$, increasing $m_0$ by one
if desired.

\paragraph{5. Low poles: cancellation of the growing number of layers.}
The negative-degree filtration of Suzuki's Proposition 2.1 works
on the root surface as well. Its proof only uses the negative
definiteness of the intersection matrix and the nonnegativity of
off-diagonal intersections; here the normalized matrix is $B/N$.
Thus $J^{-pm}/\mathcal O$ and $J^{-pm}/J^{-m}$ admit filtrations
by line bundles of strictly negative degree on the $\mathfrak C_i$.
Tensoring with $V$ gives no sections, because $V|_{\mathfrak C_i}$
is semistable of degree zero. Equation (F8.9) consequently injects
$H^0(U_mR^1\Psi M)$ into the kernel of a Frobenius-linear map minus
an injective linear map on
$H^1(J^{-m}V/V)$. The fixed-space argument from F5 bounds its
$\mathbf F_p$-dimension by $h^1(J^{-m}V/V)$.

We bound the latter dimension uniformly for the fixed $m=m_0$.
Choose the negative filtration of $J^{-m}/\mathcal O$ by adding
components to the cycle $m\mathfrak Z_N$. For a layer $P$ on
$\mathfrak C_i$ we have
\[
 -mL_i\leq\deg P<0.
 \tag{F8.12}
\]
Decompose $V|_{\mathfrak C_i}=\bigoplus_{\chi}V_{i,\chi}$ by its
generic $\mu_N$-characters, and let $r_{i,\chi}$ be their ranks.
Only the summand whose character cancels that of $P$ contributes
to cohomology. As the coefficient of $\mathfrak C_i$ increases
from $0$ to $mNc_i$, the generic character of the successive
layers runs through every residue modulo $N$ exactly $mc_i$
times. Contributions from other components have trivial generic
character at $\mathfrak C_i$. Therefore the sum of the ranks of
all contributing summands for this component is exactly
\[
 mc_i\sum_\chi r_{i,\chi}=mc_i r.
 \tag{F8.13}
\]
This is the needed cancellation of the extra $N$ layers.

Each contributing bundle $V_{i,\chi}\otimes P$ is semistable of
negative slope $\deg P$ and descends after killing the generic
gerbe. By (F8.4) and (F8.12), its first cohomology has dimension
at most its rank times $g_i+\nu_i+mL_i+1$. Summing the exact
filtration and using (F8.13) gives
\[
 h^1(J^{-m}V/V)
 \leq r\,m\sum_i c_i(g_i+\nu_i+mL_i+1).
 \tag{F8.14}
\]
The right side is independent of $N$. This bounds
$\dim_{\mathbf F_p}H^0(R^1\Psi M)$ by a constant times $r$.

\paragraph{6. The proper exceptional term.}
Let $\mathfrak D_N=\bigcup_i\mathfrak C_i$ with its reduced
structure. Its normalization sequence has one term at each crossing,
with stabilizer $\mu_N^2$. Invariants there have dimension at most
$r$ and no higher coherent cohomology. On a component, only the
generic-character-zero summand of $V$ contributes; it is semistable
of degree zero. Estimate (F8.5) thus yields
\[
 h^1(\mathfrak D_N,V)
 \leq r\left(\sum_i(g_i+\nu_i)+\#\{\text{crossings}\}\right).
 \tag{F8.15}
\]
The Artin--Schreier sequence identifies
$H^1(\mathfrak D_N,M)$ with the Frobenius-fixed subgroup of this
finite-dimensional coherent cohomology. Surjectivity of $F_V-1$
on global sections follows from the same finite-dimensional
semilinear argument as in F5. Its $\mathbf F_p$-dimension is
therefore at most (F8.15).

The nearby-cycle five-term sequence, justified by coarse descent
in step 2, combines (F8.14) and (F8.15) to prove (F8.0). One possible
constant is
\[
 C_X=m_0\sum_i c_i(g_i+\nu_i+m_0L_i+1)
            +\sum_i(g_i+\nu_i)+\#\{\text{crossings}\}.
\]
It depends on $X,J,p$ and not on the coefficient rank, the monodromy
group, or any tame inertia order.
\end{proof}

\begin{corollary}[F8b: all pro-$p$-by-pro-prime-to-$p$ quotients; Proved]
Every quotient $G$ of $\pi_1(U)$ admitting an
exact sequence $1\to P\to G\to\Gamma\to1$ with $P$ pro-$p$ and
$\Gamma$ pro-prime-to-$p$ is uniformly topologically finitely
generated. No abelian hypothesis on $\Gamma$ is necessary.
\end{corollary}
\begin{proof}
Every finite simple $\mathbf F_p[G]$-module kills $P$ and hence has
prime-to-$p$ monodromy. Such a system is tamely ramified on the
resolution. F8 and inflation give the uniform degree-one bound.
The quotient $\Gamma$ has at most a fixed number of generators,
because it is a quotient of $\pi_1(U)^{(p')}$, which is finitely
generated by Grothendieck--Murre. The relative criterion F7a gives
the conclusion. Alternating wild and tame normal factors remain
outside this assertion, as do general nonsolvable wild kernels.
\end{proof}
